\section{Model conventions and inputs}\label{sec:inputs}

Throughout the proof $q\in[1,4)$ is fixed. Constants may depend on $q$,
on a compact set of charge parameters, and on the stated geometric
separations. No estimate is asserted uniformly as $q\uparrow4$.
Write
\[
 m=\sqrt q=2\cos\lambda,\qquad
 \rho=\frac{\lambda}{\pi}\in(0,1/3],\qquad
 \kappa=\frac{4\pi}{\pi-\lambda}.
\]
We use unit square tiles in the intermediate arguments. Their corners
alternate between primal and dual sites, and the primal edges are tile
diagonals. Thus the primal mesh is $\sqrt2$ times the tile-edge length.
Each tile has one of the two noncrossing pairings of its four side
midpoints. An actual primal or dual path uses open edges of that color;
boundary identifications are not path edges. This distinction is essential
when a crossing must reach a wall or separate two arm clusters.

The two arcs in each tile can be drawn disjointly, with their order at
successive ports retained. This rounded drawing and the usual medial
drawing have uniform distance $O(1)$ in tile units under corresponding
parametrizations. In physical coordinates their distance is $O(n^{-1})$.
When a strand visits a medial vertex twice, the two occurrences retain
their traversal order. For a simply connected domain, the usual
Dobrushin completion pairs the boundary ports along the wired and free
arcs. More generally, a positive cap completion wires each designated
primal boundary interval separately and adds no primal identification
on a dual interval. The switch law with weight $m$ per completed loop
is then exactly the critical FK law, up to a deterministic factor.
Indeed, Euler's formula gives
\[
 \#\{\text{completed loops}\}=2k_\xi(\omega)+|\omega|+C,
 \qquad m^2=q,\qquad \frac{p_q}{1-p_q}=m.
\]
Consequently all conditional laws obtained by exposing switches have
the FK domain Markov property with their induced boundary partitions.
These conventions are those of \cite[Section ``Tiles, positive laws,
and localization,'' Lemma ``The positive loop law'']{FKCompanion}.

\subsection{Trace and local FK estimates}

\begin{standinginput}[Chordal trace convergence]\label{inp:trace}
For the marked uniform polygon approximations of a bounded marked
Jordan domain used in Theorem~\ref{thm:main}, the critical Dobrushin
interface converges in law to chordal $\mathrm{SLE}_\kappa$ in the
uniform Euclidean oriented-curve topology modulo increasing
reparametrization.
\end{standinginput}

This is precisely \cite[Theorem ``Dobrushin convergence for
$1\le q<4$,'' label \texttt{chordal:trace}]{FKCompanion}. We use it
as the trace input, without deriving any assertion about occupation
measures from trace convergence alone.

Here and below, a buffer has fixed positive relative width and at least
a sufficiently large fixed number of mesh layers. A collar is a full
annulus, or one of the fixed rectangular or straight-wall sector
geometries explicitly described below. The graph and deterministic
wire portions are fixed inside the collar; only the partition transmitted
from the unobserved side varies. Only the number of prescribed
deterministic through-wire blocks is bounded, and no additional
identification jumps the collar. The exterior-induced partition on
the unobserved side is arbitrary: its number of blocks is not bounded.
Nor is any bound imposed on the number of random open components
crossing the collar. The full annuli used for inward comparison have
no prescribed through-wires.

\begin{standinginput}[FK localization]\label{inp:fk}
The following estimates hold for critical square-lattice FK at the
fixed parameter $q$.
\begin{enumerate}[label=(\roman*),leftmargin=*]
\item An actual crossing of a discrete simple-polygon quadrilateral
with bounded extremal distance has probability bounded below uniformly
in the boundary partition. The dual statement holds as well. Arm
separation, quasi-multiplicativity, and localization to fixed separated
landing intervals hold for each fixed color word, in particular for
one arm and for a primal--dual pair.
\item If two positive laws agree on a collar between radii $r$ and
$16r$, their marginals beyond its fixed side have mutually bounded
densities. Equivalently, expectations of every nonnegative function
on the remote edge set are comparable. This holds for a full annulus,
for a bounded number of sectors separated by straight coordinate walls
with matching deterministic wire portions, and for a fixed-shape
rectangular bulk collar. Coincident banks of a slit are treated as
distinct boundary incidences.
\item Through $k$ sufficiently separated full annuli, the ratio of
remote expectations improves to $1+O((1-c)^k)$, or both expectations
are zero. The assertion holds in either radial direction and after
positive modifications supported on the unobserved side.
\item If $X$ counts disjoint raw strand passages through a fixed
buffered annulus or one of the preceding sectors, then, for every
fixed $B<\infty$,
\[
 \E\bigl[e^{BX}\mid\text{exterior}\bigr]\le C_B.
\]
The bound is uniform in scale and in the exterior-induced boundary
partition. It also holds for a fixed finite family of such counts.
\item Let $r_j$ be dyadic radii and $N_j$ the number of raw nests wholly
inside $r_j$, in a geometry extending beyond the largest cutoff by a
buffer. A bulk nest is a closed loop surrounding the marked point.
The increment satisfies $0\le N_j-N_{j-1}\le X_j$, where the buffer
of $X_j$ is disjoint from $B_{r_{j-k_0}}$ for a fixed $k_0$.
For $b>0$, set $F_j=\E b^{N_j}$ and
$\dd\nu_j=F_j^{-1}b^{N_j}\dd\Pp$. Then
\[
 C_b^{-1}\le\frac{F_j}{F_{j-1}}\le C_b,
 \qquad \E_{\nu_j}e^{BX_i}\le C_{b,B}\quad(i\ge j).
\]
These statements remain valid conditionally on exterior data of the
specified geometry, and with a fixed family of buffered passage counts
adjoined at or beyond the outer tilt level. The constants can be chosen
uniformly for $b$ in a compact subset of $(0,\infty)$.
\end{enumerate}
\end{standinginput}

Part~(i) is \cite[Theorem~1.2 and Propositions~6.2, 6.3, and~6.5]{DMT}.
The last three propositions apply to fixed color words, including the
two words needed here. Parts~(ii)--(v) are respectively the lemmas
``Bounded densities across a collar,'' ``Relative comparison,''
``Exponential passage moments,'' and ``Local tilt bounds'' in
\cite[Section ``Tiles, positive laws, and localization'']{FKCompanion};
their source labels are \texttt{prelim:collar},
\texttt{prelim:relative-collar}, \texttt{prelim:passages}, and
\texttt{prelim:tilts}. The compact-uniform assertion in Part~(v)
follows from those proofs by bounding $|\log b|$ on the compact set.
We require Part~(iii) only in full annuli. In particular, these inputs
do not assert comparison across arbitrary irregular walls or under
complex-valued tilts.

The separation argument with high probability used in
Section~\ref{sec:coupling} is the actual-cluster construction of
\cite[Lemma~5.9]{CDCH}, whose original statement concerns FK-Ising.
For our range of $q$, its conditional crossing steps are supplied by
\cite[Theorem~1.2]{DMT}, as in the proof of that paper's
Proposition~6.2. Section~\ref{sec:coupling} specifies the extra buffers
and the arm continuations needed for the coupling; we do not identify
ordinary separation up to constants with that stronger conclusion.

The free and wired plane laws coincide, are invariant under lattice
symmetries, and are exchanged by primal--dual symmetry at these
self-dual weights; see \cite[Theorem~1.1]{DCLM} and its square-lattice
specialization in \cite[Section ``Crossing estimates and finite plane
loops'']{FKCompanion}. Buffered circuits of both colors show that
every plane medial strand is almost surely a finite loop. A translation
by one tile may interchange the two colors, but preserves this loop law.

\subsection{Plane loops and angular rows}

\begin{standinginput}[The nested plane loop limit]\label{inp:loops}
The entire loop collection of the free infinite-volume critical
square-lattice FK law converges, through the full mesh sequence, to
nested whole-plane $\mathrm{CLE}_\kappa$. Convergence is in spherical
loop-matching topology, with curves compared modulo their traversal
parametrizations, and includes all nested generations. The limit is
invariant under similarities and, more generally, M\"obius maps.
\end{standinginput}

This is the separately cited Theorem ``Nested full-plane convergence
for $q<4$,'' label \texttt{loop:plane}, in
\cite[Section ``From squares to the free plane law and the
sphere'']{FKCompanion}. Its statement concerns exactly the plane law
used here. It is stronger than Input~\ref{inp:trace}, and is used for
the finite-resolution nesting limits in Sections~\ref{sec:radial}
and~\ref{sec:calibration}. Passage moments from Input~\ref{inp:fk}
provide the required uniform integrability when nesting functionals
are unbounded.

Orienting a completed contractible loop and assigning total turning
weights $e^{i\lambda}$ and $e^{-i\lambda}$ recovers its unoriented
weight $m$. Summing the local orientations gives the homogeneous
six-vertex weights
\[
 a=b=1,\qquad c=2\cos(\lambda/2),\qquad
 \Delta=\frac{2-c^2}{2}=-\cos\lambda.
\]
The induced height changes by one across an arrow. The plane arrow
law is the zero-slope law obtained by first taking the longitudinal
limit of balanced even-width tori and then letting the width tend to
infinity. Here balance means zero total arrow flux across a transverse
section. The identification of finite local arrow and phase words
with the FK loop expansion is part of the plane-kernel construction
in \cite[Section ``Row algebra and the infinite-plane state'']{FKCompanion}.
It agrees with the balanced-plane law in \cite[Theorem~2.2 and
Equation~(18)]{DKLM}.

\begin{standinginput}[Physical angular representation]\label{inp:rows}
At $0<\lambda\le\pi/3$, the normalized homogeneous rows have a plane
Hilbert-space representation with vacuum $\Omega$. Finite ordinary
spin and diagonal phase words have the expectations of the preceding
physical plane law. With the angular coordinate normalized so that
the strip is $|\Rea z|<\pi/2$, marked rows $A^D(z)$ are bounded and
holomorphic, satisfy $\|A^D(z)\|\le2\|D\|$, and obey
\[
 A^D(z+ib)=U_bA^D(z)U_b^{-1},\qquad U_b=e^{ibB},
\]
where $B$ is self-adjoint. For a diagonal unitary twist, the normalized
commuting row family has scalar spectral functions which, under
$w=e^{iz}$, are inner in the right half-plane. They continue without
zeros through a common-width band at each strip edge. In particular,
their zeros lie in a fixed cone $\Rea w_j\ge c_*|w_j|$, and their
general factorization has the form
\[
 f(w)=\eta\exp(-a_*w-b_*/w)\prod_j\beta_j(w),\qquad
 a_*,b_*\ge0,\quad |\eta|=1,
\]
where $\beta_j$ are normalized half-plane Blaschke factors and
$\sum_j\Rea w_j/(1+|w_j|^2)<\infty$. The sum may initially be infinite.
\end{standinginput}

These are the plane-kernel, angle-generator, and fixed-band
continuation lemmas, labels \texttt{bulk:plane-kernels},
\texttt{bulk:boost}, and \texttt{bulk:edge-band}, together with
Equation \texttt{bulk:inner-factorization}, in
\cite[Section ``Angular transfer and its two leading modes'']{FKCompanion}.
The argument in \cite[Section~4]{GFFCompanion} for the current
cocycle, finite-inner spectrum, spectral-width tails, and variogram
is stated there for $\pi/3<\lambda<\pi$. Those stronger conclusions
are not included in Input~\ref{inp:rows}. In
Section~\ref{sec:spectral} we transfer the needed arguments to the
representation above and prove the conclusions at our parameter.

\subsection{The covariance normalization}

Only a separated-increment covariance, rather than a pointwise
Gaussian limit for the height, is needed to determine the electric
constant. Let $h_\delta$ be the height of the preceding balanced plane
law on mesh $\delta$, with an arbitrary additive constant, and put
$G(u,v)=-(2\pi)^{-1}\log|u-v|$.

\begin{standinginput}[Separated height covariance]\label{inp:covariance}
For four distinct points $u,u',v,v'$ and lattice approximants tending
to them,
\begin{align*}
 &\lim_{\delta\downarrow0}
 \E\bigl[(h_\delta(u')-h_\delta(u))
          (h_\delta(v')-h_\delta(v))\bigr]\\
 &\qquad=\sigma^2\bigl(G(u',v')-G(u',v)-G(u,v')+G(u,v)\bigr),
 \qquad \sigma^2=\frac{2}{\pi-\lambda}.
\end{align*}
The convergence is locally uniform away from endpoint collisions.
\end{standinginput}

Indeed, $c=2\cos(\lambda/2)\in[\sqrt3,2)$, so the covariance
part of \cite[Definition~2.7 and Theorem~2.8]{DKLM} applies directly
to this physical law and unit-height convention. Its multiplier is
$1/\arcsin(c/2)=2/(\pi-\lambda)$. The same normalization is recorded
in the main theorem of \cite{GFFCompanion}. In
Section~\ref{sec:electric} we use Input~\ref{inp:covariance} only
for four distinct collinear points. Thus this use needs neither the
smaller-$c$ part of that companion's GFF theorem nor an extension
of the DKLM theorem beyond its stated range.
